\section{合同的变更和转让}
\begin{enumerate}
  \item 必考，结合保全，1 个多选题。
  \item 生效要件 + 法律效力；这 2 个必须要签协议。
  \item \zd{箭头指向债权人，为未来的债权转让、债务承担、债权人代位权之诉、债权人撤销权之诉的行权能结合在一起，捋的会比较清楚}。
  \item 债权转让/准物权变动。换人要债。\zd{内外有别 +（双重）通知}。\al{银行处理坏账；保理合同}。
        \begin{enumerate}
          \item \al{乙欠甲 30W，丁提供房屋抵押并登记。甲欠乙 20W，甲转让 30W 给丙}。
          \item 生效要件。合法有效的债权及债权转让协议、具有可让与性。
          \item 对债务人和担保人的影响都不大。
          \item 协议通知生效。
                \begin{enumerate}
                  \item 新的民法典，谁通知都可以了；新债权人通知要有证据。
                  \item 方式随意。
                  \item 通知后原债权人后悔不能撤销，需新债权人同意。
                \end{enumerate}
          \item \zd{债权转让担保的从属性}。\zd{担保自动转让，从随主}。\al{丁的房屋抵押转让给丙}。
          \item 履行费用。谁产生谁付。
          \item 诉讼时效中断，重新起算 3 年 \yjt 债的同一性原理。\al{乙可主张诉讼时效抗辩}。
          \item 让与人对其让与的债权负\zd{瑕疵担保责任 \yjt 违约责任}。
          \item \zd{抵销权}。互债、同类、主动债权已到期。\zd{到期在先原则}。谁先到期保护谁。
        \end{enumerate}
  \item 债务承担/债务转移。内容不变主体变。\zd{先判断是退出还是加入机制}。
        \begin{enumerate}
          \item 生效要件。合法有效的债务及债务承担协议、具有可让与性。
          \item 类型。
                \begin{enumerate}
                  \item 退出机制/免责的债务承担。债务承担协议 +\zd{（双重顺序）同意}（沉默视为不同意。没考过）。
                  \item 加入机制/并存的债务承担。\zd{协议 + 通知}。后果。连带责任，有约从约，无约全部。新增考点。
                \end{enumerate}
        \end{enumerate}
  \item 第三人代为履行/清偿。\zd{无协议 + 主动清偿}了全部 \hz 部分债务。
\end{enumerate}
